
% === Append to merged_exercises_FULL_v2.tex ===
\subsection*{Exercise 1.2 (Periodic Green's function for $-\psi''+\psi=f$)}
Suppose $f\in C^0(\mathbb{R})$ is $2\pi$-periodic, i.e.\ $f(\theta)=f(\theta+2\pi)$. For $\theta\in[0,2\pi)$ define
\begin{equation}\label{eq:psi-split}
\psi(\theta):=\int_{0}^{\theta} f(\alpha)\,\frac{\cosh(\pi-\theta+\alpha)}{2\sinh\pi}\,d\alpha
+\int_{\theta}^{2\pi} f(\alpha)\,\frac{\cosh(-\pi-\theta+\alpha)}{2\sinh\pi}\,d\alpha,
\end{equation}
and extend $\psi$ to $\mathbb{R}$ by periodicity $\psi(\theta+2\pi)=\psi(\theta)$.

\paragraph{Background / definition.}
It is convenient to write \eqref{eq:psi-split} using the periodic Green kernel
\begin{equation}\label{eq:Gdef}
G(\theta,\alpha):=\frac{\cosh\!\big(\pi-|\theta-\alpha|\big)}{2\sinh\pi},\qquad \theta,\alpha\in[0,2\pi).
\end{equation}
Then \eqref{eq:psi-split} is equivalently
\begin{equation}\label{eq:psi-conv}
\psi(\theta)=\int_{0}^{2\pi} f(\alpha)\,G(\theta,\alpha)\,d\alpha,
\end{equation}
i.e.\ $\psi=L^{-1}f$ for the periodic operator $L:=-\frac{d^2}{d\theta^2}+1$.

\begin{solution}
\textbf{(a) $\boldsymbol{\psi\in C^0(\mathbb{R})}$.}
On $[0,2\pi)$ the map $(\theta,\alpha)\mapsto f(\alpha)\,G(\theta,\alpha)$ is continuous and bounded.
By dominated convergence, $\theta\mapsto\psi(\theta)$ is continuous. The periodic extension is continuous
since both $f$ and $G$ are $2\pi$-periodic in each variable.

\medskip
\textbf{(b) First derivative and $\boldsymbol{\psi\in C^1(\mathbb{R})}$.}
Differentiate \eqref{eq:psi-split} by Leibniz' rule. Using
\(
\partial_\theta\cosh(\pi-\theta+\alpha)=-\sinh(\pi-\theta+\alpha)
\)
and
\(
\partial_\theta\cosh(-\pi-\theta+\alpha)=-\sinh(-\pi-\theta+\alpha),
\)
and noting that the boundary terms cancel because
$
\frac{\cosh\pi}{2\sinh\pi}-\frac{\cosh\pi}{2\sinh\pi}=0,
$
we obtain
\begin{equation}\label{eq:psi-prime}
\psi'(\theta)= -\int_0^\theta f(\alpha)\,\frac{\sinh(\pi-\theta+\alpha)}{2\sinh\pi}\,d\alpha
             \;-\; \int_\theta^{2\pi} f(\alpha)\,\frac{\sinh(-\pi-\theta+\alpha)}{2\sinh\pi}\,d\alpha.
\end{equation}
Again by dominated convergence, the right-hand side is continuous, hence $\psi\in C^1(\mathbb{R})$.

\medskip
\textbf{(c) Second derivative and the equation $\boldsymbol{-\psi''+\psi=f}$.}
Differentiate \eqref{eq:psi-prime} once more. Using
\(
\partial_\theta\sinh(\pi-\theta+\alpha)=-\cosh(\pi-\theta+\alpha)
\)
and
\(
\partial_\theta\sinh(-\pi-\theta+\alpha)=-\cosh(-\pi-\theta+\alpha),
\)
Leibniz' rule yields
\begin{align*}
\psi''(\theta)
&=-\frac{f(\theta)\sinh\pi}{2\sinh\pi}+\int_{0}^{\theta} f(\alpha)\,\frac{\cosh(\pi-\theta+\alpha)}{2\sinh\pi}\,d\alpha \\
&\quad-\frac{f(\theta)\sinh(-\pi)}{2\sinh\pi}+\int_{\theta}^{2\pi} f(\alpha)\,\frac{\cosh(-\pi-\theta+\alpha)}{2\sinh\pi}\,d\alpha \\
&=-f(\theta)+\psi(\theta).
\end{align*}
Thus for all $\theta\in\mathbb{R}$, we have
\begin{equation*}
-\psi''(\theta)+\psi(\theta)=f(\theta).
\end{equation*}
Since $f$ is continuous, the display shows $\psi\in C^2(\mathbb{R})$.
\end{solution}

\paragraph{Figures (for intuition).}
Place the PNGs \texttt{kernel\_heatmap.png}, \texttt{f\_and\_psi.png}, and \texttt{residual\_check.png}
next to your \LaTeX{} file, then include:
\begin{figure}[h]
  \centering
  \includegraphics[width=.75\linewidth]{kernel_heatmap.png}
  \caption{Green's kernel $G(\theta,\alpha)=\frac{\cosh(\pi-|\theta-\alpha|)}{2\sinh\pi}$.}
\end{figure}
\begin{figure}[h]
  \centering
  \includegraphics[width=.75\linewidth]{f_and_psi.png}
  \caption{Example input $f$ and the corresponding solution $\psi=L^{-1}f$.}
\end{figure}
\begin{figure}[h]
  \centering
  \includegraphics[width=.75\linewidth]{residual_check.png}
  \caption{Numerical residual $-\psi''+\psi-f$ (near $0$).}
\end{figure}
